\chapter{Detalles técnicos da implementación}
\label{app:detalles-implementacion}

Este anexo amplía os aspectos técnicos descritos no
Capítulo~\ref{chap:implem}. Inclúense exemplos das interfaces proporcionadas
pola \gls{api}, fragmentos representativos do código e detalles internos da
implementación que non resultan necesarios para comprender o funcionamento
xeral da solución, pero permiten documentar con maior precisión as decisións
adoptadas.

\section{Detalles da implementación da API}

\subsection{Documentación e endpoints dispoñibles}

FastAPI xera automaticamente unha especificación OpenAPI a partir das rutas,
dos modelos de entrada e saída e das respostas declaradas. Esta especificación
permite consultar e executar os endpoints desde unha interface Swagger sen
necesidade dun cliente adicional.

% IMAXE 1
% Facer unha captura da pantalla inicial de Swagger na que aparezan:
% Procesamento de informes
% Preprocesamento de informes
% Centros sanitarios
% Sistema
% Comunicación P2P
\IfFileExists{imaxes/anexos/implementacion/swagger-api.png}{
\begin{figure}[htbp]
    \centering
    \includegraphics[width=0.95\textwidth]
    {imaxes/anexos/implementacion/swagger-api.png}
    \caption{Documentación OpenAPI xerada automaticamente pola API.}
    \label{fig:anexo-swagger}
\end{figure}
}{}

Na Figura~\ref{fig:anexo-swagger} poden observarse as operacións agrupadas
segundo a súa responsabilidade. O procesamento principal realízase mediante
\textit{/extraccion/}, mentres que o resto das rutas permiten comprobar o
estado do servizo, consultar centros sanitarios, probar o preprocesamento de
forma independente e reenviar as mensaxes destinadas á sincronización entre
dispositivos.

\begin{table}[htbp]
    \centering
    \caption{Principais endpoints implementados na API.}
    \label{tab:anexo-endpoints}
    \resizebox{\columnwidth}{!}{
    \begin{tabular}{p{0.16\textwidth} p{0.27\textwidth} p{0.47\textwidth}}
        \hline
        \textbf{Método} & \textbf{Endpoint} & \textbf{Finalidade} \\
        \hline
        POST & /extraccion/ &
        Procesar unha folla e devolver os datos estruturados. \\

        POST & /preprocesamento/ &
        Aplicar unicamente o preprocesamento e devolver o ficheiro resultante. \\

        GET & /centros/buscar &
        Recuperar a información dun centro sanitario a partir do seu nome. \\

        GET & /health/ &
        Comprobar o estado e a versión da API. \\

        POST & /notificar/enviar &
        Reenviar mediante Firebase Cloud Messaging unha mensaxe preparada polo cliente. \\
        \hline
    \end{tabular}
    }
\end{table}


\subsection{Exemplo de extracción}

A Figura~\ref{fig:anexo-extraccion} mostra unha chamada realizada directamente
desde Swagger. A entrada da operación é o documento orixinal, mentres que a
resposta xa contén a información transformada á estrutura utilizada polo
cliente móbil.

% IMAXE 2
% Captura de POST /extraccion/
% Ideal: que se vexa o ficheiro seleccionado, código 200 e o inicio do JSON.
% Usa unha folla de proba, nunca datos persoais reais.
\IfFileExists{imaxes/anexos/implementacion/chamada-extraccion.png}{
\begin{figure}[htbp]
    \centering
    \includegraphics[width=0.95\textwidth]
    {imaxes/anexos/implementacion/chamada-extraccion.png}
    \caption{Exemplo dunha extracción realizada desde Swagger.}
    \label{fig:anexo-extraccion}
\end{figure}
}{}

Neste exemplo pode apreciarse que o cliente non recibe a saída orixinal de
Azure Document Intelligence. A API realiza previamente a interpretación e
validación dos campos e constrúe unha resposta formada pola cabeceira, o
calendario de doses, o histórico e os metadatos da análise.

Unha resposta simplificada, empregando datos ficticios, presenta a seguinte
estrutura:

\begin{lstlisting}[
    caption={Exemplo simplificado da resposta de /extraccion/.},
    label={lst:json-extraccion}
]
{
  "cabeceira": {
    "dataInforme": "2025-01-14",
    "inr": "2,5",
    "farmaco": "Sintrom 4 mg",
    "doseSemanal": "13,5 mg",
    "proximaVisita": "21/01/2025",
    "centro": "C.S. EXEMPLO"
  },
  "calendario": [
    {
      "data": "2025-01-15",
      "dia": 15,
      "dose": "1/2",
      "accion": "TOMAR",
      "eControl": false,
      "diaSemanaTexto": "MERCORES"
    }
  ],
  "historico": [],
  "metadatos": {
    "confianzaGlobal": 0.94,
    "modelo": "M1"
  }
}
\end{lstlisting}


\subsection{Modelos de resposta}

O contrato devolto ao cliente defínese mediante modelos Pydantic. Estes modelos
permiten establecer os campos obrigatorios, os tipos esperados e a estrutura
xerárquica da resposta antes de que FastAPI a serialice como \gls{json}.

O seguinte fragmento mostra dous dos modelos empregados para representar o
calendario e a resposta completa:

\begin{lstlisting}[
    language=Python,
    caption={Fragmento dos modelos Pydantic empregados pola API.},
    label={lst:pydantic-response}
]
class DoseDia(BaseModel):
    data: str
    dia: int
    dose: str | None
    accion: str
    eControl: bool
    diaSemanaTexto: str

class AnalisisResponse(BaseModel):
    cabeceira: CabeceiraResponse
    calendario: List[DoseDia]
    historico: List[ItemHistorico]
    metadatos: MetadatosResponse
\end{lstlisting}

A separación en modelos específicos evita construír unha resposta sen unha
estrutura definida e serve ademais como base para a documentación OpenAPI.


\subsection{Inicialización dos servizos externos}

Durante o ciclo de vida de FastAPI créase o cliente de Azure Document
Intelligence unha única vez e almacénase no estado da aplicación. As rutas
poden recuperalo posteriormente mediante inxección de dependencias.

\begin{lstlisting}[
    language=Python,
    caption={Inicialización e recuperación do cliente de Azure.},
    label={lst:azure-dependency}
]
client = DocumentIntelligenceClient(
    endpoint=DOC_INTEL_ENDPOINT,
    credential=AzureKeyCredential(api_key),
    api_version="2024-11-30"
)
app.state.doc_intel_client = client

# ...

def get_azure_client(request: Request):
    client = request.app.state.doc_intel_client
    if not client:
        raise HTTPException(
            status_code=500,
            detail="Servizo Azure non dispoñible"
        )
    return client
\end{lstlisting}

Esta implementación permite reutilizar a mesma conexión durante a vida da API
e facilita tamén a súa substitución por unha implementación simulada durante
as probas.


\subsection{Preprocesamento das imaxes}

O preprocesamento implementouse de maneira conservadora. O seu obxectivo non é
transformar todas as imaxes recibidas, senón corrixilas unicamente cando
existen evidencias suficientes dunha inclinación ou deformación da folla.

Inicialmente compróbase o tipo e o tamaño do documento. Os PDF e os formatos
que non poden ser modificados de maneira segura continúan sen cambios. Nas
imaxes decodificables redúcese temporalmente a resolución empregada para a
detección, convértese a imaxe a escala de grises e extráense os bordos.

A detección tenta localizar primeiro un contorno cuadrilateral compatible coa
xeometría dunha folla. Se esta estratexia non proporciona un resultado fiable,
téntase localizar unha superficie clara e, posteriormente, estimar a
inclinación a partir das liñas detectadas. A corrección só se aplica cando se
superan os limiares de confianza establecidos. Ante un resultado dubidoso ou
un erro, consérvase o ficheiro orixinal.

O seguinte fragmento resume esta estratexia:

\begin{lstlisting}[
    language=Python,
    caption={Selección conservadora do preprocesamento dunha imaxe.},
    label={lst:preprocesamento}
]
candidato = _buscar_folla(gris, bordos)
if candidato is not None:
    cuadrilatero, confianza = candidato
    resultado = _corrixir_perspectiva(
        contido,
        imaxe,
        cuadrilatero / escala,
        confianza,
        tipo,
    )
    if resultado is not None:
        return resultado

candidato_claro = _buscar_folla_clara(
    imaxe_deteccion,
    gris,
    bordos,
)

# Se non existe unha correccion fiable,
# conservase o documento orixinal.
\end{lstlisting}

% IMAXE 3
% A mellor figura do anexo.
% Esquerda: foto inclinada tomada co móbil.
% Dereita: saída de /preprocesamento/.
% Emprega a MESMA imaxe nas dúas partes.
\IfFileExists{imaxes/anexos/implementacion/preprocesamento-comparativa.png}{
\begin{figure}[htbp]
    \centering
    \includegraphics[width=0.95\textwidth]
    {imaxes/anexos/implementacion/preprocesamento-comparativa.png}
    \caption{Comparación entre o documento orixinal e o resultado do preprocesamento.}
    \label{fig:anexo-preprocesamento}
\end{figure}
}{}

Na Figura~\ref{fig:anexo-preprocesamento} obsérvase o efecto da corrección sobre
unha fotografía tomada con perspectiva. O documento resultante presenta unha
xeometría máis próxima á dunha folla escaneada, facilitando a posterior análise
polo modelo documental. Se a detección non alcanza unha confianza suficiente,
esta transformación non se aplica.

Para facilitar a comprobación desta etapa implementouse ademais
\textit{/preprocesamento/}. A resposta inclúe cabeceiras de diagnóstico que
indican se se aplicou unha transformación, o método empregado, o motivo, o
ángulo detectado, a confianza e as dimensións de entrada e saída.


\subsection{Endpoints auxiliares}

O endpoint \textit{/health/} proporciona unha comprobación mínima da
dispoñibilidade da API:

\begin{lstlisting}[
    caption={Exemplo de resposta do endpoint de estado.},
    label={lst:health-json}
]
{
  "status": "OK",
  "timestamp": "2026-08-27T20:15:32.425",
  "version": "1.0.0"
}
\end{lstlisting}

Para os centros sanitarios utilízase un catálogo KML. Antes da procura
normalízase o nome recibido, eliminando diferenzas de maiúsculas, acentos e
algunhas variacións na forma de escribir ``C.S.''. Os candidatos compáranse
posteriormente mediante similitude textual.

A resposta presenta unha estrutura como a seguinte:

\begin{lstlisting}[
    caption={Exemplo ilustrativo da resposta dunha procura de centro.},
    label={lst:centro-json}
]
{
  "nome": "CENTRO DE SAUDE EXEMPLO",
  "enderezo": "Rua Exemplo, 1",
  "concello": "Exemplo",
  "telefono": "982000000",
  "fonte": "Xunta de Galicia - Mapa dos servizos sanitarios (CC BY-SA 4.0)"
}
\end{lstlisting}

Finalmente, \textit{/notificar/enviar} non interpreta o contido clínico
intercambiado entre dispositivos. Recibe un token, un tipo de aviso e o
\textit{payload} xa cifrado polo dispositivo de orixe e utiliza Firebase Cloud
Messaging para reenvialo. Unha resposta correcta contén unicamente o resultado
do envío e o identificador xerado por Firebase:

\begin{lstlisting}[
    caption={Estrutura da resposta dun envío mediante Firebase.},
    label={lst:notificacion-response}
]
{
  "success": true,
  "message_id": "projects/.../messages/..."
}
\end{lstlisting}



\section{Detalles da implementación do cliente móbil}

\subsection{Inicialización e recuperación do estado}

O proceso de inicialización concentra as operacións que deben completarse antes
de mostrar a interface. \textit{main} prepara Flutter e chama a
\textit{InicializadorAplicacion}, que inicializa Firebase e as notificacións,
rexistra os receptores de Firebase Cloud Messaging e recupera a configuración
persistente.

Unha vez coñecidos o estado da configuración e o rol gardado, a aplicación
selecciona a ruta inicial. Se se trata dun paciente previamente configurado,
tamén se reconstrúe a planificación dos recordatorios da pauta activa.

% IMAXE 4
% Eu FARIA ESTA EN DRAW.IO, non captura.
% Participantes:
% main
% InicializadorAplicacion
% Firebase
% Flutter Secure Storage
% LocalNotificationService
% DatabaseService
% MyApp
%
% Secuencia:
% main -> Inicializador: inicializar()
% Inicializador -> Firebase: initializeApp()
% Inicializador -> Firebase: onBackgroundMessage / onMessage
% Inicializador -> Notificacions: inicializar()
% Inicializador -> SecureStorage: configuracion_finalizada + rol_usuario
% [Paciente] -> SecureStorage: hora_toma
% [Paciente] -> DatabaseService: obterEstados()
% [Paciente] -> Notificacions: programarTomasPaciente()
% Inicializador --> main: ResultadoInicializacion
% main -> MyApp: runApp()
\IfFileExists{imaxes/anexos/implementacion/inicializacion-cliente.png}{
\begin{figure}[htbp]
    \centering
    \includegraphics[width=0.92\textwidth]
    {imaxes/anexos/implementacion/inicializacion-cliente.png}
    \caption{Secuencia de inicialización do cliente móbil.}
    \label{fig:anexo-inicializacion}
\end{figure}
}{}

A Figura~\ref{fig:anexo-inicializacion} permite observar que a selección da
pantalla inicial non depende unicamente dunha variable en memoria. O estado
persistente recupérase antes de construír a aplicación e, no caso do paciente,
tamén se comproba a planificación das notificacións.


\subsection{Estado da captura e revisión}

O modelo de vista encargado da captura mantén un estado reducido composto polo
nome do ficheiro, a posible mensaxe de erro e un indicador booleano que sinala
se existe unha operación en curso:

\begin{lstlisting}[
    caption={Estado mantido por ReportCaptureViewModel.},
    label={lst:report-capture-state}
]
bool _procesando = false;
String? _nomeFicheiro;
String? _erro;

bool get procesando => _procesando;
String? get nomeFicheiro => _nomeFicheiro;
String? get erro => _erro;
\end{lstlisting}

A vista observa estes valores mediante \textit{Provider}. Cando
\textit{procesando} é verdadeiro, as opcións de selección substitúense por un
indicador de progreso. O propio modelo de vista non contén elementos gráficos,
senón unicamente o estado empregado pola interface.

A revisión posterior permite actualizar ou eliminar días e modificar a
seguinte visita. Antes de confirmar a pauta execútanse as seguintes
comprobacións:

\begin{itemize}
    \item debe existir polo menos unha entrada no calendario;
    \item as datas deben ser válidas;
    \item unha mesma data non pode aparecer dúas veces;
    \item os días normais deben conter unha dose;
    \item admítense doses numéricas e fraccións cun denominador diferente de cero;
    \item os días de control non requiren dose;
    \item se existe unha data de seguinte visita, esta tamén debe ser válida.
\end{itemize}

% IMAXE 5 OPCIONAL
% Captura da pantalla de revisión cun erro provocado expresamente:
% data duplicada ou dose non válida.
% Ao lado, mesma pauta xa corrixida.
\IfFileExists{imaxes/anexos/implementacion/validacion-pauta.png}{
\begin{figure}[htbp]
    \centering
    \includegraphics[width=0.88\textwidth]
    {imaxes/anexos/implementacion/validacion-pauta.png}
    \caption{Exemplo da detección e corrección dun dato durante a revisión da pauta.}
    \label{fig:anexo-validacion}
\end{figure}
}{}

A Figura~\ref{fig:anexo-validacion} ilustra a función da revisión manual como
última comprobación antes de converter os datos extraídos nun tratamento
activo.


\subsection{Estrutura da persistencia local}

SQLite mantén catro conxuntos principais de información:

\begin{table}[htbp]
    \centering
    \caption{Táboas creadas pola aplicación móbil.}
    \label{tab:sqlite-anexo}
    \resizebox{\columnwidth}{!}{
    \begin{tabular}{p{0.25\textwidth} p{0.66\textwidth}}
        \hline
        \textbf{Táboa} & \textbf{Información almacenada} \\
        \hline
        pautas &
        Data, día, dose, acción, control, día da semana, estado,
        hora de confirmación e desviación en minutos. \\

        analise\_actual &
        Data do informe, seguinte visita, centro, fármaco, dose semanal,
        INR e identificador do documento. \\

        historico\_visitas &
        Visitas anteriores, INR, fármaco, doses, seguinte visita e comentarios. \\

        pacientes\_coidador &
        UID, token, nome, resumo sincronizado e momento da última actualización
        das persoas supervisadas. \\
        \hline
    \end{tabular}
    }
\end{table}

O gardado completo dunha análise execútase mediante unha transacción. Antes de
substituír a pauta actual calcúlase un identificador a partir dos campos
principais do documento e do calendario. Se coincide co identificador xa
almacenado, recupéranse os estados das datas existentes.

\begin{lstlisting}[
    caption={Recuperación dos estados ao importar novamente a mesma folla.},
    label={lst:mesma-folla}
]
final documentoId = _identificadorDocumento(analise);

final eMesmaFolla =
    actual.isNotEmpty &&
    actual.first['documentoId'] == documentoId;

final rexistrosAnteriores =
    <String, Map<String, Object?>>{};

if (eMesmaFolla) {
  final filas = await txn.query(
    'pautas',
    columns: [
      'data',
      'estado',
      'horaConfirmacion',
      'desviacionMinutos'
    ],
  );

  for (final fila in filas) {
    rexistrosAnteriores[fila['data'] as String] = fila;
  }
}
\end{lstlisting}

Deste xeito, a reimportación dunha folla xa coñecida non elimina a información
de seguimento rexistrada previamente.


\subsection{Estados dunha toma}

Cada día con dose efectiva pode pasar polos seguintes estados:

\begin{table}[htbp]
    \centering
    \caption{Estados empregados para o seguimento das tomas.}
    \label{tab:estados-toma}
    \begin{tabular}{p{0.31\textwidth} p{0.58\textwidth}}
        \hline
        \textbf{Estado} & \textbf{Significado} \\
        \hline
        PENDENTE &
        A toma aínda non foi confirmada. \\

        TOMADA &
        A persoa confirmou a toma sen superar a hora prevista. \\

        TOMADA\_FORA\_HORA &
        A confirmación produciuse despois da hora configurada. \\

        NON\_TOMADA &
        O día rematou sen que a toma fose confirmada. \\
        \hline
    \end{tabular}
\end{table}

Cando se confirma unha toma calcúlase a diferenza en minutos respecto da hora
habitual e almacénase xunto co instante de confirmación. Isto permite conservar
non só se a dose foi tomada, senón tamén cando se realizou.


\subsection{Planificación das notificacións}

A creación dos recordatorios realízase independentemente da API concreta de
notificacións. \textit{PlanificadorRecordatorios} recibe a pauta, a hora
configurada e os estados e devolve os avisos que aínda deben producirse.

Para unha toma pendente poden existir dous recordatorios: un á hora
configurada e un segundo, por defecto trinta minutos máis tarde, se continúa
sen confirmar. Os estados xa pechados non xeran novos avisos.

\begin{lstlisting}[
    caption={Cálculo dos dous recordatorios asociados a unha toma.},
    label={lst:recordatorios}
]
final toma = DateTime(
  ano,
  mes,
  dia,
  horaNumero,
  minutoNumero,
);

final esquecemento = toma.add(
  Duration(minutes: marxeEsquecementoMinutos),
);

if (toma.isAfter(agora)) {
  // Engadir aviso da toma
}

if (esquecemento.isAfter(agora)) {
  // Engadir aviso de esquecemento
}
\end{lstlisting}

% IMAXE 6
% Facer dúas capturas reais de Android:
% 1. "Hora da toma"
% 2. "Toma sen confirmar"
% Podes xuntalas nunha soa figura.
\IfFileExists{imaxes/anexos/implementacion/notificacions.png}{
\begin{figure}[htbp]
    \centering
    \includegraphics[width=0.75\textwidth]
    {imaxes/anexos/implementacion/notificacions.png}
    \caption{Notificación da hora da toma e aviso posterior de toma sen confirmar.}
    \label{fig:anexo-notificacions}
\end{figure}
}{}

A Figura~\ref{fig:anexo-notificacions} mostra os dous tipos de aviso xerados
localmente. Se a toma se confirma antes do segundo aviso, a notificación de
esquecemento é cancelada.


\subsection{Vinculación segura entre dispositivos}

A vinculación comeza no dispositivo do paciente. Para cada tentativa xérase un
identificador novo e unha clave provisional de 256 bits. O código QR contén a
versión do protocolo, o identificador da vinculación, o UID e o token Firebase
do paciente e a clave provisional.

A seguinte representación emprega datos deliberadamente ficticios:

\begin{lstlisting}[
    caption={Estrutura do código QR empregado durante a vinculación.},
    label={lst:qr-vinculacion}
]
{
  "version": 1,
  "vinculacionId": "K8x...ficticio",
  "uidPaciente": "uid-de-proba",
  "tokenPaciente": "token-fcm-de-proba",
  "clave": "clave-base64-de-proba"
}
\end{lstlisting}

% IMAXE 7
% Composición:
% esquerda: QR no paciente
% dereita: cámara do supervisor escaneándoo
% O QR DEBE SER DE PROBA.
\IfFileExists{imaxes/anexos/implementacion/vinculacion-qr.png}{
\begin{figure}[htbp]
    \centering
    \includegraphics[width=0.90\textwidth]
    {imaxes/anexos/implementacion/vinculacion-qr.png}
    \caption{Inicio da vinculación entre paciente e persoa supervisora mediante código QR.}
    \label{fig:anexo-qr}
\end{figure}
}{}

Na Figura~\ref{fig:anexo-qr} represéntase a única etapa na que é necesaria unha
acción presencial entre os dous dispositivos. O QR proporciona os datos
iniciais necesarios para establecer unha canle segura, pero a clave incluída
nel non é a empregada para as comunicacións posteriores.


\subsection{Intercambio e almacenamento das claves}

Tras ler o QR, o dispositivo supervisor crea unha nova clave permanente de
256 bits. Esta clave envíase ao paciente dentro da primeira mensaxe, protexida
coa clave provisional do QR.

O proceso pode resumirse nas seguintes etapas:

\begin{enumerate}
    \item o paciente xera o identificador e a clave provisional;
    \item o supervisor escanea o QR;
    \item o supervisor xera unha nova clave permanente;
    \item a clave permanente envíase cifrada coa clave provisional;
    \item o paciente autentica e descifra a primeira mensaxe;
    \item ambos dispositivos conservan a clave permanente mediante
    \textit{Flutter Secure Storage};
    \item a clave provisional pendente elimínase no paciente.
\end{enumerate}

Deste xeito, unha captura posterior do código QR non proporciona a clave que se
está empregando nas comunicacións dunha vinculación xa establecida.


\subsection{Formato das mensaxes cifradas}

As mensaxes posteriores empréganse con AES-GCM de 256 bits. Para cada envío
xérase un \textit{nonce} novo e o resultado serialízase nun envelope coa
seguinte estrutura:

\begin{lstlisting}[
    caption={Estrutura dun envelope cifrado AES-GCM.},
    label={lst:aes-envelope}
]
{
  "version": 1,
  "vinculacionId": "...",
  "nonce": "...",
  "datos": "...",
  "mac": "..."
}
\end{lstlisting}

O identificador da vinculación e o tipo de aviso non se inclúen no texto
cifrado, senón que se empregan tamén como datos autenticados adicionais:

\begin{lstlisting}[
    caption={Construción dos datos autenticados adicionais.},
    label={lst:aad}
]
utf8.encode(
  'sintrom-p2p-v1|$id|$tipoAviso'
);
\end{lstlisting}

Isto vincula criptograficamente o contido ao seu contexto. Se se modifica o
tipo de mensaxe, o identificador, o texto cifrado ou o código de autenticación,
o descifrado non pode completarse correctamente.


\subsection{Tipos de mensaxes sincronizadas}

A sincronización utiliza o campo \textit{tipo\_aviso} para indicar a operación
que debe executar o receptor.

\begin{table}[htbp]
    \centering
    \caption{Principais mensaxes empregadas na sincronización P2P.}
    \label{tab:tipos-p2p}
    \resizebox{\columnwidth}{!}{
    \begin{tabular}{p{0.37\textwidth} p{0.52\textwidth}}
        \hline
        \textbf{Tipo} & \textbf{Finalidade} \\
        \hline
        VINCULACION\_INICIAL &
        Completar o intercambio inicial de claves. \\

        ESTADO\_COMPLETO &
        Transferir ao supervisor o estado actual do paciente. \\

        SOLICITAR\_SINCRONIZACION &
        Solicitar novamente o estado actualizado. \\

        TOMA\_CONFIRMADA &
        Informar da confirmación dunha dose. \\

        TOMA\_ESQUECIDA &
        Informar dunha toma que quedou sen confirmar. \\

        NOVO\_INFORME &
        Indicar que a pauta do paciente foi modificada. \\

        CONFIGURACION\_ACTUALIZADA &
        Sincronizar cambios de nome ou hora habitual. \\

        INFORME\_FRAGMENTO &
        Transferir unha parte dunha folla completa. \\

        DESVINCULAR\_PACIENTE /
        DESVINCULAR\_COIDADOR &
        Eliminar unha vinculación nos dous extremos. \\
        \hline
    \end{tabular}
    }
\end{table}


\subsection{Fragmentación dunha folla completa}

Unha análise completa pode ter un tamaño elevado para enviala directamente
nunha única mensaxe. Por este motivo, \textit{enviarInformeRemoto} serializa a
análise e divídea en fragmentos de ata 2200 caracteres.

Cada fragmento contén un identificador común, a súa posición dentro da
secuencia, o número total de partes e a porción correspondente dos datos:

\begin{lstlisting}[
    caption={Estrutura empregada para fragmentar un informe.},
    label={lst:fragmento-informe}
]
{
  "id": id,
  "indice": i,
  "total": total,
  "datos": json.substring(i * tamanho, fin)
}
\end{lstlisting}

No receptor, cada parte almacénase temporalmente empregando unha clave formada
polo identificador do informe e o índice do fragmento. A reconstrución só
continúa cando están dispoñibles todas as partes.

\begin{lstlisting}[
    caption={Reconstrución dun informe a partir dos fragmentos recibidos.},
    label={lst:reconstrucion-informe}
]
final partes = <String>[];

for (var i = 0; i < total; i++) {
  final parte = await _storage.read(
    key: 'informe_${id}_$i',
  );

  if (parte == null) return;
  partes.add(parte);
}

final analise = AnaliseModel.fromJson(
  jsonDecode(partes.join()),
);

await _db.gardarAnalise(analise);
\end{lstlisting}

% IMAXE 8 OPCIONAL
% Diagrama moi pequeno:
%
% AnaliseModel
%     |
% serializacion JSON
%     |
% +--------+--------+--------+
% | frag 0 | frag 1 | frag 2 |
% +--------+--------+--------+
%     |        |        |
%   AES-GCM  AES-GCM  AES-GCM
%     \        |        /
%        dispositivo paciente
%                |
%         reconstruccion
%                |
%          AnaliseModel
%
\IfFileExists{imaxes/anexos/implementacion/fragmentacion-informe.png}{
\begin{figure}[htbp]
    \centering
    \includegraphics[width=0.82\textwidth]
    {imaxes/anexos/implementacion/fragmentacion-informe.png}
    \caption{Fragmentación, transferencia e reconstrución dunha folla enviada desde o dispositivo supervisor.}
    \label{fig:anexo-fragmentacion}
\end{figure}
}{}

A Figura~\ref{fig:anexo-fragmentacion} resume este proceso. Cada fragmento
mantén a protección criptográfica da comunicación e o tratamento activo só se
actualiza cando se puido reconstruír correctamente a análise completa. Unha vez
gardada, reprográmanse os recordatorios e elimínanse os fragmentos temporais.